% Figure 1.1: the six design principles (problem, rule, where it is enforced)
\begin{figure}[H]
  \centering
  % \principle{icon}{title}{problem}{rule}{where}
  \newcommand\principle[5]{%
    \begin{tcolorbox}[card]
      \cardtitle{#1}{#2}
      {\color{black!55}\itshape #3}\par
      #4
      \tcblower
      \faArrowRight\ #5
    \end{tcolorbox}}
  \begin{tcbraster}[raster columns=3, cardgrid, raster before skip=0pt]
    \principle{\faProjectDiagram}{A constrained flow}{Models are not predictable.}
      {The model interprets; code decides which \acsp{api} to call.}{\S\ref{sec:routing}}
    \principle{\faUserLock}{Identity outside the model}{Prompts can be manipulated.}
      {The customer ID comes only from the login token (\ac{jwt}), never from a prompt.}{\S\ref{sec:state-context}}
    \principle{\faHistory}{The right edition}{Old editions still apply.}
      {Clauses filtered by product and signed edition.}{\S\ref{sec:knowledge-base}}
    \principle{\faCalculator}{Numbers from the APIs}{Wrong numbers cost most.}
      {Amounts and dates come from the core system and are checked.}
      {\texttt{output\_guard}, Figure~\ref{fig:graph}}
    \principle{\faEyeSlash}{No unverified text}{Drafts can be wrong.}
      {Only progress states until the output guardrail passes.}{\S\ref{sec:api}}
    \principle{\faEuroSign}{Zero cost, minimal exposure}{Free services see what they get.}
      {Only what the answer needs leaves the server, identifiers masked.}
      {\S\ref{sec:deployment}}
  \end{tcbraster}
  \caption{The six design principles.}
  \label{fig:principles}
\end{figure}
